\documentclass[a4paper, 10pt, openany]{report}


% \documentclass[draft, a4paper, 10pt, openany]{book}
\usepackage[a4paper]{geometry}

\usepackage[spanish, es-lcroman]{babel}

\usepackage{parskip}
\usepackage{multicol}

\usepackage{subfiles} 
\usepackage{tocbibind}
\usepackage[toc]{appendix}

\usepackage[justification=centering]{caption}
\usepackage{subcaption}

\usepackage[shortlabels]{enumitem}
\usepackage{csquotes}
\usepackage{epigraph}

\usepackage{fontspec}
% switch to a monospace font supporting more Unicode characters
\setmonofont{JetBrains Mono}[Scale=MatchLowercase]
\usepackage[outputdir=../build]{minted}
\newmintinline[lean]{lean4}{bgcolor=white}
% \newminted[leancode]{lean4}{fontsize=\footnotesize}
\usemintedstyle{tango}  % a nice, colorful theme


\usepackage{booktabs}

% \usepackage[Conny]{fncychap}

\usepackage{../sty/teoremas}
\usepackage{../sty/general}
\usepackage{../sty/conjuntos}
\usepackage{../sty/calculo}
\usepackage{../sty/algebra}

% Se carga el último 
\usepackage{bookmark}

% Los apéndices se llaman apéndices
\renewcommand\appendixpagename{Apéndices}
\renewcommand{\appendixtocname}{Apéndices}

% mainmatter no reinicia el contador de páginas
% \makeatletter
% \renewcommand\mainmatter{
%     \@mainmattertrue\cleardoublepage\renewcommand\thepage{\arabic{page}}}
% \makeatother

\renewcommand{\thesubsection}{\thesection.\alph{subsection}}

\providecommand{\keywords}[1]
{
  \small	
  \textbf{\textit{Palabras clave---}} #1
}


\title{Resultados clásicos sobre semántica de lenguajes de programación en Lean 4}
\author{Ricardo Maurizio Paul}
% \date{2020--2024}

\begin{document}

% \frontmatter{}

\maketitle

% \clearpage
%\cleardoublepage % openright
% \vspace*{\fill}
% \thispagestyle{empty} % Sin número de página
% \begin{quotation}
% 	\raggedleft{}
% 	\em % énfasis
% 	Duc, o parens celsique dominator poli, \\
% 	Quocumque placuit; nulla parendi mora est. \\
% 	Adsum inpiger. Fac nolle, comitabor gemens \\
% 	Maiusque patiar, facere quod licuit bono. \\
% 	Ducunt volentem fata, nolentem trahunt. \\
% 	\bigskip
% 	--- Seneca
% \end{quotation}
% \vspace*{\fill}

\tableofcontents

\begin{abstract}
	Desarrollamos algunos resultados clásicos sobre la semántica de lenguajes de 
	programación utilizando el asistente de demostración Lean 4. 
	Definimos un pequeño lenguaje imperativo \texttt{IMP} junto a
	sus semánticas de paso largo, paso corto y denotacional.
	Analizamos algunas de sus propiedades y demostramos su equivalencia.

	\keywords{
		asistente de demostración,
		Lean 4, 
		métodos formales, 
		semántica, 
	}
\end{abstract}

% Es necesario que las etiquetas de los capítulos estén en mayúsculas para que
% puedan ser referenciadas en el header de los apéndices :/

\chapter{Introducción}\label{CHAP:INTRO}
\subfile{intro.tex}

\chapter{Lean}\label{CHAP:LEAN}
% \epigraph{
%   A las 3 se va pal' beauty, a las 4 va pal' mall \\
%   Pa' eso de las 5 está a vapor \\
%   Se hizo las boobies y el booty, viste de Christian Dior \\
%   Activa a toas' sus amigas con un call
% }{Lean --- Superiority}
\subfile{lean.tex}

\chapter{Teoría de tipos vs teoría de conjuntos}\label{CHAP:SETS}
\subfile{sets.tex}

\chapter{El lenguaje \texttt{IMP}}\label{CHAP:IMP}
\subfile{imp.tex}


% \part{Apéndices}
% \appendix

% \chapter{Problemas resueltos}

% \section{Problemas del capítulo~\ref{CHAP:EC_LINEAL_PRIMER_ORDEN}}
% \subfile{2_ec_lin_primer_orden/problemas}

% \section{Problemas del capítulo~\ref{CHAP:SISTEMAS}}
% \subfile{3_sistemas_lin/problemas}

% \chapter{Resumen de diagramas de fase}
% \subfile{4_cualitativo/anexo}

% \backmatter{}

% \listoffigures

% \listoftables

\bibliographystyle{alpha}
\bibliography{bibliografia}
% \nocite{*}

\end{document}
